\documentclass[10pt,a4paper]{amsart}
\usepackage[margin=19mm]{geometry}
\usepackage{amsmath,amssymb,mathtools}
\usepackage[scr=rsfs,cal=euler]{mathalfa}
\usepackage{xfrac}
\usepackage[unicode,hidelinks,bookmarks=false]{hyperref}
\newcommand{\ie}{i.e.\ }
\newcommand{\cf}{cf.\ }
\newcommand{\resp}{respectively\ }
\newcommand{\Subsection}{\subsection}
\providecommand{\nobreakdash}{\nobreak\hbox{-}}
\setlength{\emergencystretch}{2em}
\multlinegap0pt
\usepackage{bm}
\usepackage{bbm}
\usepackage{dsfont}
\usepackage{xspace}
\usepackage{cancel}
\usepackage{mathtools}
\usepackage{amsthm}
\makeatletter
\@ifundefined{theorem}{\newtheorem{theorem}{Theorem}}{}
\@ifundefined{proposition}{\newtheorem{proposition}[theorem]{Proposition}}{}
\@ifundefined{sublemma}{\newtheorem{sublemma}{}}{}
\makeatother
\newcommand{\cl}{\text{cl}}
\title[Weak maps and the Tutte Polynomial]{Weak maps and the Tutte Polynomial}
\author{Christine Cho, James Oxley}
\date{Source: arXiv:2501.15606v3 (CC BY 4.0)}
\begin{document}
\maketitle

\noindent\emph{Source and licence.} Weak maps and the Tutte Polynomial. Version arXiv:2501.15606v3 (CC BY 4.0), distributed under the Creative Commons Attribution 4.0 International licence, as recorded in the arXiv licence metadata for this exact version and shown on the abstract page https://arxiv.org/abs/2501.15606v3. Full rights notice: licenses/arxiv-2501.15606-CC-BY-4.0.txt.\par\smallskip

\noindent\textit{Editorial scope.} Counted unit: the Proposition whose source label is \texttt{hypclones} in the article (source0.tex lines 490--492, source label \texttt{hypclones}), delivered together with the article's own complete original proof of it (source0.tex lines 494--553, one \texttt{proof} environment of 60 lines). No mathematics is written, repaired or re-derived: statement and proof are the article's own text line for line. Required context retained uncounted: none. Every definition, display and label of those lines is retained unchanged. Omitted from this package and not redistributed: 1--489, 493--493, 554--626. Nothing inside a retained excerpt is omitted or altered: every retained body is delivered exactly as the article's lines are written, so no omission receipt of any kind accompanies this package. Citations: the retained excerpts carry no citation command at all, so no bibliography entry of the article is transcribed and no reference is redistributed. Reference adaptation: none was needed and none was made; every reference inside the delivered unit resolves inside it, so no unresolved reference or citation is produced. Printed slips kept literal: no printed word, symbol, hypothesis or number of the counted statement or of its proof was altered. Layout and class: the article's own class is \texttt{amsart} with packages including enumitem, graphicx, tikz, tkz-graph, wrapfig, float, caption, subcaption, amsmath, amssymb, pgfplots, xcolor, harpoon; the wrapper is the lane's pinned \texttt{amsart} editorial wrapper at 10pt on A4, which restates the theorem-like environments the retained text needs and adds the article's own macro definitions that the retained text needs (\texttt{\textbackslash cl}), with extra wrapper packages (bm, bbm, dsfont, xspace, cancel, mathtools). The delivered numbering is the unit's own. Assets: no retained span contains a figure environment, a graphics-inclusion command, a PDF-page inclusion, a table or a TikZ picture; no image file is copied into this package, the package declares no asset dependency and no image file is redistributed with it. Licence: version arXiv:2501.15606v3 of this article is distributed under the Creative Commons Attribution 4.0 International licence, as recorded in the arXiv licence metadata for this exact version and shown on the abstract page https://arxiv.org/abs/2501.15606v3. A full rights notice is delivered at licenses/arxiv-2501.15606-CC-BY-4.0.txt. Characters: every retained excerpt is pure ASCII, so the delivered engine, XeLaTeX with its default Latin Modern text font, reports no missing character.
\begin{proposition}\label{hypclones}
    Let $f$ be freer than $g$ in $M$. Let $L$ be obtained from $M$ by deleting every element of $E(M)-\{f,g\}$ that is parallel to $g$. Then $h(f;M)=h(g;M)$ if and only if $f$ and $g$ are clones in $L$. 
\end{proposition}

\begin{proof}
    Suppose $f$ and $g$ are clones in $L$.
    As $h(f;M)=h(f;L)$ and $h(g;M)=h(g;L)$, we have $h(f;M)=h(g;M)$. 
    To prove the converse, suppose that $h(f;M)=h(g;M)$. 
    Let $\mathcal{H}(M;f;\overline g)$ be the set of hyperplanes of $M$ containing $f$ but not $g$. 
    Then $|\mathcal{H}(M;f;\overline g)|=|\mathcal{H}(M;g;\overline f)|$.
    Hence $|\mathcal{H}(L;f;\overline g)|=|\mathcal{H}(L;g;\overline f)|$.
    Clearly, if $\{f,g\}$ is a $2$-circuit of $L$, then $f$ and $g$ are clones in $L$. 
    Thus, we may assume that no 2-circuit of $L$ contains $g$.


    For $J\in \mathcal{H}(L;g;\overline f)$, let $B_J$ be an arbitrarily chosen basis of $J$ containing $g$. 
    Then $\cl_{L}(B_J-g)$ is a rank-$(r-2)$ flat of $L$. 
    Let $\cl_{L}(B_J-g)\cup f=F$. 
    Then $r(F)=r-1$ otherwise $f\in \cl_{L}(B_J-g)$, so $g\in \cl_{L}(B_J-g)$, a contradiction. 
    Consider $\cl_{L}(F)$. 
    Assume $f$ is not a coloop of $\cl_{L}(F)$. 
    Then $\cl_{L}(F)$ contains a circuit $C$ containing $f$. 
    Since $f$ is freer than $g$, we see that $g\in \cl_{L}(F)$. 
    Then $\cl_{L}(F)\supseteq B_J$. 
    Thus $\cl_{L}(F)\supseteq J$. 
    As $r(\cl_{L}(F))=r(J)$, we deduce that $\cl_{L}(F)=J$. 
    But $f\not\in J$, a contradiction. 
    Thus $f$ is a coloop of $\cl_{L}(F)$, so $\cl_{L}(B_J-g)\cup f\in \mathcal{H}(L;f;\overline g)$. 
    Let $\psi(J)=\cl_{L}(B_J-g)\cup f$. 
    Note that $\psi$ depends upon the choices made for the bases $B_J$. Moreover, $\psi$ maps $\mathcal{H}(L;g;\overline f)$ to $\mathcal{H}(L;f;\overline g )$.  

    \begin{sublemma}
       $\psi$ is bijective. 
    \end{sublemma}

    To see that $\psi$ is injective, suppose that, for distinct members $J_1$ and $J_2$ of $\mathcal{H}(L;g;\overline f)$, the hyperplanes $\psi(J_1)$ and $\psi(J_2)$ are equal. 
    Then $\cl_{L}(B_{J_1}-g)=\cl_{L}(B_{J_2}-g)$. 
    Now the rank-$(r-2)$ flat $J_1\cap J_2$ contains $g$ and so contains the rank-$(r-1)$ set $B_{J_1}$, a contradiction. 
    Since $|\mathcal{H}(L;g;\overline f)|=|\mathcal{H}(L;f;\overline g)|$ and $\psi$ is injective, we conclude that (12.1) holds.


    \begin{sublemma}
        $g$ is a coloop of $L|J$ for every $J\in \mathcal{H}(L;g;\overline f)$.  
    \end{sublemma}

    Suppose $g$ is not a coloop of $L|J$. 
    Then $\cl_{L}(B_J-g)$ is a subset of $J$ avoiding $g$. 
    As $g$ is not a coloop, there is an element $h$ of $J-\cl_{L}(B_J-g)-g$. 
    Since no 2-circuit of $L$ contains $g$, the elements $g$ and $h$ are not parallel. 
    Thus $\{g,h\}$ is independent. 
    Extend $\{g,h\}$ to a basis $B_J'$ of $L|J$. 
    Then $\cl_{L}(B_J'-g)\neq\cl_{L}(B_J-g)$ because $h\in \cl_{L}(B_J'-g)-\cl_{L}(B_J-g)$. 
    Thus $\cl_{L}(B_J'-g)\cup f$ is a member of $\mathcal{H}(L;f;\overline g)$ that is not in the set $\psi(\mathcal{H}(L;g;\overline f))$.
    As $|\mathcal{H}(L;g;\overline f)|=|\mathcal{H}(L;f;\overline g)|$ and $\psi$ is a bijection, this is a contradiction. Thus (12.2) holds. 

    \smallskip

    Suppose $g$ is not freer than $f$ in $L$.
    Then $L$ has a cyclic flat $K$ containing $g$ and avoiding $f$. 
    Take a basis $B$ for $K$ and consider $B \cup f$. 
    Extend $B\cup f$ to get a basis $B_{L}$ for $L$. 
    Then $\cl_L(B_{L}-f)$ is a hyperplane of $L$ containing $g$ and avoiding $f$. 
    Moreover, since $K\subseteq \cl_L(B_{L}-f)$, the hyperplane $\cl_L(B_{L}-f)$ has a circuit containing $g$; that is, $g$ is not a coloop of $\cl_L(B_{L}-f)$, a contradiction to (12.2).
       \end{proof} 

\end{document}
